\scanpage{60}
\mspage{66}

\deleted{Axiome suivant} \deleted{\uline{Der $5'$} (g) : \inserted{$\unsure{\mathbb D|\underline D}$ dér. à g.} Soit $f:X\longrightarrow Y$ tel que $f$ soit}

\uline{Proposition 4} \inserted{$\mathbb D$ dér. à \unsure{g \& d.|dr.|g et dr.}} Soit $f:X\longrightarrow Y$ un foncteur de $\mathrm{Diag}$, tel que pour $y\in Y$, $ {}_{y\backslash}X$ soit $\mathbb D$-contractile (donc $f$ une $\mathbb D$-équivalence costricte, sous réserve que $\mathbb D$ soit dérivateur à \unsure{droit|droite} \ldots) Alors pour \uline{tout} $F$ de $\mathbb D(Y)$ (pas néc. à coeff. constants) on a

\[ H^\bullet_{\mathbb D}(Y,F)\xrightarrow{\sim}H^\bullet_{\mathbb D}(X,f^*(F)) \]

est un iso. \unsure{Réciproquem$^{\mathrm t}$|Réciproquement} si la conclusion est vraie pour tout $F$, alors \deleted{$f$ est une} les $ {}_{y\backslash}X$ sont $\mathbb D$-contractiles.

Dém. \deleted{\unsure{Il}} revient au $=$ de dire que pour tout $M\in\mathrm{Ob}\,A=\mathbb D(e)$, on a

\[ \mathrm{Hom}(M,H^\bullet_{\mathbb D}(Y,F))\xrightarrow{\mathrm{iso}}\mathrm{Hom}(M,H^\bullet_{\mathbb D}(Y,\unsure{f_{\scriptscriptstyle -}|F|\cdot})) \]

\diagram{\begin{tikzcd}[ampersand replacement=\&,row sep=tiny,column sep=small]
X \arrow[rr,"f"] \arrow[dr,dash,"p"'] \& {} \& Y \arrow[dl,dash,"q"] \\
{} \& e \& {}
\end{tikzcd}}

Mais par adjonction, cette flèche s'interprète comme

\[ \mathrm{Hom}(M_Y,F)\longrightarrow\mathrm{Hom}(M_X,f^*F)\tag*{($*$)} \]

D'autre part, \deleted{$M_X=f^*(M_Y)$, \unsure{d'où}}, si $f_!$ existe, on a

\[ \mathrm{Hom}(M_X,f^*F)=\mathrm{Hom}(\deleted{\unsure{M}}f_!(M_X),F), \]

et l'hom ($*$) est déduit de l'hom d'adjonction

\[ f_!(\underbrace{\deleted{\unsure{M}}f^*(M_Y)}_{M_X})\longrightarrow M_Y,\tag*{($**$)} \]

en prenant les Hom dans $F$. Donc la conclusion équivaut à dire que cet hom est iso, quel que soit \deleted{\unsure{$M$}} le \deleted{\unsure{$F$}} coeff constant $M_Y$ sur $Y$. Pour voir si c'est un iso, on regarde fibre par fibre (Der 2), -- or les fibres s'interprètent comme $H_\bullet^{\mathbb D}({}_{y\backslash}X,M_X)$ \inserted{(par \uline{Der 4} (d))} et comme \deleted{\unsure{$y$}}$\unsure{M_{Y,y}|M_y}$.
% archive sheet number: 59
